\subsection{Gobierno, verificación, formación y servicio}
Las cuentas restantes distinguen el control del proyecto 1.1, las pruebas integrales 1.10.2, la formación 1.11.2/3, el cierre de marcha blanca 1.11.5 y el servicio/salida 1.12.3--11. Se instancian 317 paquetes fijos, con 7.504 HH: 2.528 de implementación y 4.976 de operación. Las fichas comunes de trabajos periódicos se aplican a cada código/ventana del registro \path{componentes/paquetes_gobierno_servicio.csv}; no se suma otra vez el componente padre.

El expediente quincenal de gobierno incluye actas, tablero, acuerdos, riesgos, cambios y valor ganado con SPI/CPI. La reversibilidad se presenta dentro de los primeros 90 días, se actualiza y se prepara para una transferencia acompañada al menos 90 días reales. La salida se programa en 53--56 para reservar más de tres meses y absorber la duración real; el calendario verificará fechas efectivas. La transferencia contractual no se confunde con reconstruir por completo una plataforma alternativa estimada en SD4.

Las pruebas integrales incluyen QA/cobertura, regresión, carga/estrés, operación desconectada, resiliencia, conmutación real de DR, seguridad ofensiva, accesibilidad y UAT/T-17. Antes de cada paso a producción, incluidas liberaciones de Operación, se comprobarán carga/estrés a 1,5 veces el peak declarado, WCAG 2.2 AA documentada, ofensiva independiente sin críticos/altos abiertos, resiliencia y DR real conforme a T-10/PRUEBAS-6. Los dos ensayos por etapa permanecen en 1.10.1 y no se repiten en sus costos. En operación, DR/resiliencia se repiten además al menos semestralmente y ofensiva al menos anualmente, con fechas relativas sujetas a ventanas permitidas y sin exceder la frecuencia comprometida. Proyecto verifica esfuerzo y revisión disponibles antes de seleccionar una liberación; si las bolsas previstas no bastan, Synaptix aporta refuerzo a su costo sin reducir las pruebas \parencite[cap.~20, secc.~20.1, pp.~34--35]{bases-transversales}.

\subsection{Sesiones de capacitación por perfil, sitio y cohorte}
Se distinguen cinco perfiles: usuarios finales, avanzados, administradores, responsables técnicos designados por el cliente y soporte niveles 1/2. Synaptix no presume un área TI interna. Los materiales incluyen fuentes editables en español, guías, FAQ, guion/video y base de conocimiento; su disponibilidad en autoformación registra avance. Se combinarán sesiones presenciales por sitio, sincrónicas y acompañamiento de tarea. Los perfiles administradores/técnicos deberán evaluarse y certificarse antes de cada cierre \parencite[RT-22.01--07]{bases-transversales}.

\paragraph{1.11.7.e.p.s.b: sesión y evaluación de cohorte}
Por etapa $e$, perfil $p$, sitio/turno $s$ y lote $b$ de hasta doce asistentes, el entregable incluye preparación, dos sesiones de dos horas, ejercicio y registro individual. José Lara Arce responde por el paquete: Implantación 8, Funcional 2 y Calidad 2 HH, total 12. La aceptación requiere ejercicio reproducido, asistencia, resultado y recuperación asignada; asistencia sin evaluación no acredita certificación. Se requieren material recibido, entorno seguro, usuarios autorizados y ventana acordada. Excluye participación del cliente, traslado, formación incremental de innovación y sesiones de recuperación adicionales no realizadas.

El número de paquetes será $B_e=\sum_{p,s}\lceil N_{e,p,s}/12\rceil$, sobre personas y turnos confirmados, sin asumir que el mismo material reduce automáticamente el número de sesiones. Se agregan $12(B_1+B_2)$ HH de implementación. En operación se usará \texttt{1.12.10.m.p.s.b}, con mes $m$, para incorporación de personal nuevo o recuperación requerida, con la misma ficha y costo incluidos en servicio. Las dos jornadas anuales fijas no sustituyen esa incorporación. La coordinación de personal estacional y monitores nuevos debe completar formación durante noviembre y antes del 15 de diciembre, sin capacitación dentro del peak protegido. Cero participantes sólo se registrará para una cohorte confirmada sin personas.

\subsection{Presencia por ola y estabilización}
\paragraph{1.11.4.e.w.d.z: acompañamiento diario en terreno}
Cada etapa $e$ y ola $w$ genera paquetes por día $d$ y zona $z$. Se propone un puesto de acompañamiento de ocho horas por zona durante los primeros catorce días y cuatro horas por zona en los siguientes catorce, en ventanas seguras, con disminución condicionada a adopción. El paquete diario tendrá 8 HH: las primeras dos semanas, ocho de presencia; las siguientes, cuatro de presencia y cuatro de preparación, seguimiento y evidencia. Excluye reuniones de control, capacitación formal y conciliación de datos ya estimadas. El referente de coordinación será José Lara Arce; la ejecución requerirá personal de implantación de refuerzo con competencias y turnos acreditados, sin atribuir esas jornadas al titular además de sus otros paquetes.

Se recibirá registro de zona, horario, tareas acompañadas, incidencias, pendientes y entrega al relevo. Debe existir dotación por turno, transporte y suplencia; se cubrirán fines de semana aplicables. El traslado se planifica por separado de las horas de acompañamiento. Para dos zonas por ola y 28 días son 56 paquetes y 448 HH, con 336 horas de presencia y 112 de preparación/seguimiento. Si una ola requiere otras zonas o más cobertura, se instanciarán sus paquetes adicionales.

La estabilización de cuatro semanas después de cada paso a producción empleará \texttt{1.11.6.e.d.z}, dos puestos de ocho horas diarios, uno en pantalanes y otro en varadero, durante 28 días consecutivos: 56 paquetes/448 HH por etapa, incluidos fines de semana. Se recibirá evidencia de atención, relevo y resolución. E1 pertenece a implementación; E2 al servicio de operación del mes 21. La dotación en esos puestos requiere reemplazos y horarios autorizados; dos puestos no significan que dos personas puedan trabajar 28 días sin descanso. La presencia queda separada de las 96 HH de estabilización de datos por etapa, que realizan tareas técnicas distintas.

\subsection{Cobertura externa y especialistas}
La solución reservará cobertura NOC y SOC continua y centro de atención bilingüe 06:00--24:00 en temporada alta y 08:00--20:00 el resto, con emergencia y planes vencidos 24$\times$7. Se propone un puesto activo de NOC, uno de SOC y al menos uno de atención por turno; capacidad de cola, competencias, supervisión y relevo deberán demostrar los SLA. No se presentan proveedores, ubicaciones ni contratos como seleccionados o firmados \parencite[RT-21.01/06]{bases-transversales}\parencite[RT-21.06]{caso07}.

Para un mes con $D_m$ días, cada puesto continuo requiere $24D_m$ horas de cobertura; atención requiere $18D_{H,m}+12D_{B,m}$, donde temporada alta es 15 de noviembre--31 de marzo. Los puestos continuo de NOC/SOC deben recibir alertas y escalar críticos/altos fuera del horario de atención, sin espera hasta apertura. Con 144 horas de cobertura efectiva por equivalente de jornada como hipótesis de dimensionamiento, las cotas son $\lceil24D_m/144\rceil$ por puesto continuo y $\lceil(18D_{H,m}+12D_{B,m})/144\rceil$ para atención. Ese supuesto se contrastará con turnos legales, descansos, ausencias y reemplazos del prestador; no acredita dotación contratada.

Las horas de cobertura se presupuestan como capacidad externa reservada y no se atribuyen a los ocho titulares de SD1. Su precio incluye recepción y tareas de primer nivel dentro de la cobertura contratada; esas horas no se vuelven a sumar como atención individual. El gerente dedicado y especialistas atienden coordinación, diagnóstico y resolución especializada dentro de sus paquetes y disponibilidad. Volumen de incidencias, tiempos de servicio, concurrencia y reemplazo definirán puestos adicionales para sostener 80\,\% antes de veinte segundos, resolución al primer contacto de 70\,\% y abandono máximo 5\,\%; un solo puesto no demuestra esos resultados.
